\chapter{Implementacja i ewaluacja systemu detekcji}

Celem implementacji było zbudowanie praktycznego narzędzia do wsadowej analizy \textbf{plików wideo} pod kątem obecności jawnych i ukrytych znaków wodnych, metadanych pochodzenia oraz artefaktów typowych dla materiałów wygenerowanych lub przetworzonych przez systemy generatywnej~AI. System powstał w~języku \textbf{Python}, który oferuje dojrzały ekosystem bibliotek do przetwarzania obrazu, analizy sygnału i~budowy interfejsów desktopowych, a~zarazem umożliwia szybkie prototypowanie i~iteracyjne rozszerzanie metod detekcji.

Architektura rozwiązania ma charakter modułowy. Aplikacja składa się z~trzech warstw: interfejsu użytkownika (GUI), warstwy orkiestracji analizy oraz zestawu wyspecjalizowanych detektorów odpowiedzialnych za wykrywanie watermarków widocznych, próbę odczytu watermarków niewidzialnych, analizę metadanych C2PA, ocenę artefaktów częstotliwościowych i~identyfikację statycznych nakładek powtarzających się w~kolejnych klatkach. Taki podział upraszcza dalszy rozwój projektu, ponieważ nowe heurystyki i~klasyfikatory mogą być dodawane jako osobne moduły bez przebudowy całej aplikacji.

\section{Język programowania i biblioteki}

Rdzeń systemu zaimplementowano w~języku \textbf{Python~3}. Graficzny interfejs użytkownika zbudowano z~wykorzystaniem biblioteki \textbf{PyQt6} --- zestawu wiązań języka Python do biblioteki Qt6, który pozwala tworzyć natywnie wyglądające aplikacje okienkowe na systemach Windows, macOS i~Linux. Dzięki temu możliwe było przygotowanie w~pełni funkcjonalnej aplikacji desktopowej z~tabelą wyników, mechanizmem kolejkowania plików, paskiem postępu, podglądem klatki oraz obsługą raportów generowanych dla każdego analizowanego materiału.

W warstwie przetwarzania wideo kluczową rolę pełnią biblioteki~\textbf{OpenCV} i~\textbf{NumPy}. OpenCV odpowiada za odczyt klatek z~pliku wideo, operacje na~obrazie (np. filtrowanie CLAHE, wyrównanie histogramu, detekcja krawędzi), obliczanie przepływu optycznego metodą Farnebäcka oraz wstępne przygotowanie danych dla kolejnych etapów. NumPy służy do wydajnych obliczeń macierzowych i~operacji na~pikselach. Biblioteka \textbf{Pandas} umożliwia agregację wyników z~wielu klatek i~zapis raportów do~plików~CSV.

Do rozpoznawania tekstu w~klatkach wideo (watermarki widoczne, napisy, \textit{lower thirds}) aplikacja korzysta z~silników~OCR. W~zależności od konfiguracji środowiska uruchamiany jest silnik \textbf{EasyOCR} lub \textbf{PaddleOCR} --- oba działają lokalnie i~nie wymagają połączenia z~usługami zewnętrznymi. W~warstwie analizy statystycznej sygnału wideo oraz~przy obliczaniu widma częstotliwościowego klatek (analiza FFT) wykorzystywana jest biblioteka \textbf{SciPy}.

Detekcja watermarków niewidzialnych i~sprawdzenie metadanych C2PA (\textit{Content Credentials}) realizowane są przez dedykowane moduły~\texttt{c2pa\_detector} i~dodatkowe komponenty opierające się na~publicznych specyfikacjach standardu C2PA opracowanego przez~\textit{Coalition for Content Provenance and Authenticity}. W~części badawczej uwzględniono również wpływ systemów~\textbf{SynthID} (Google DeepMind) na materiały generowane przez popularne generatory, co pozwoliło zestawić własny detektor z~rozwiązaniami opartymi na~znakach osadzanych bezpośrednio w~procesie generacji.

\section{Architektura aplikacji i GUI}

Interfejs użytkownika napisany w~PyQt6 pełni rolę warstwy sterującej całym procesem analitycznym. Główne okno aplikacji podzielone jest na~dwa panele rozdzielone ruchomym separatorem (\textit{splitter}): lewy panel zawiera tabelę wyników, konsolę logów i~pasek postępu, natomiast prawy panel prezentuje podgląd aktualnie analizowanej klatki oraz powiększony widok obszaru, w~którym wykryto watermark (\textit{Zoom}).

Tabela wyników zawiera sześć kolumn: nazwę pliku, typ materiału, status detekcji AI, procentowy udział klatek z~wykryciem, status metadanych C2PA oraz ścieżkę do raportu~CSV. Dwukrotne kliknięcie wiersza otwiera w~eksploratorze plików katalog wynikowy przypisany do~danego materiału wideo.

Obsługiwane formaty wideo to: \texttt{mp4}, \texttt{mov}, \texttt{avi}, \texttt{mkv} i~\texttt{webm}. Aplikacja jest ukierunkowana wyłącznie na~analizę \textbf{materiałów filmowych}, nie obrazów statycznych, co odpowiada celowi pracy --- detekcji watermarków i~śladów generatywnego pochodzenia w~plikach wideo, gdzie istotna jest zarówno zawartość pojedynczej klatki, jak i~informacja temporalna między klatkami.

Pliki mogą być dodawane do~kolejki przez okno dialogowe, przez przeciągnięcie i~upuszczenie (\textit{drag and drop}) lub rekurencyjne skanowanie całego folderu. Użytkownik konfiguruje czułość OCR (parametr \textit{confidence}), gęstość próbkowania klatek (\textit{sample rate}) oraz opcjonalny tryb szczegółowej analizy dwufazowej, który po~pierwszym przebiegu ponownie analizuje klatki bez detekcji z~zastosowaniem zaawansowanych filtrów obrazu (CLAHE, Top-Hat, korekta gamma, odwrócone kolory).

\section{Moduły detekcji}

Warstwa detekcji składa się z~kilku współpracujących modułów uruchamianych kolejno dla każdego pliku wideo. Moduł \texttt{ocr\_detector} odpowiada za~ekstrakcję klatek, normalizację obrazu i~wykrywanie tekstu metodami OCR. Wykryte regiony tekstowe są filtrowane według progu pewności i~porównywane ze~słownikiem znanych fraz watermarkingowych (\enquote{Generated by AI}, \enquote{Runway}, \enquote{Sora} i~inne). Wyniki są zapisywane w~raporcie wraz ze~współrzędnymi wykrytych obszarów, co umożliwia późniejsze przejrzenie klatek z~zaznaczonymi detekcjami.

Analiza statycznych nakładek polega na~obliczeniu mediany temporalnej sekwencji klatek --- element powtarzający się przez wiele kolejnych klatek w~tym~samym miejscu (np. logo stacji) daje silny sygnał w~medianie, który odróżnia go~od~treści dynamicznych. Analiza przepływu optycznego pozwala z~kolei wykryć obszary o~zerowym lub prawie zerowym ruchu, które mogą wskazywać na~obecność nakładki statycznej.

Dla każdego pliku wytwarzany jest zestaw artefaktów diagnostycznych: klatki z~zaznaczonymi obszarami detekcji, obraz mediany temporalnej, wizualizacje z~analizy FFT oraz plik~\texttt{report.csv} zawierający metadane analizy, numery klatek, współrzędne wykryć i~poziomy pewności. Taki sposób dokumentowania wyników sprawia, że system nie działa jak \enquote{czarna skrzynka} --- użytkownik może przejrzeć każdy artefakt i~zweryfikować zasadność decyzji klasyfikacyjnej.

\section{Opis ilustracji i zrzutów ekranu}

% TODO: wstawić zrzuty ekranu z folderu photos i uzupełnić opisy poniżej

Każda ilustracja w~tym rozdziale powinna opisywać \textbf{konkretny etap działania aplikacji}, a~nie wyłącznie widok interfejsu. Oznacza to, że podpis rysunku powinien wskazywać: który moduł jest aktywny, co wyświetla tabela lub panel podglądu, jakie artefakty są widoczne i~jaki jest sens prezentowanego wyniku.

Przykładowe schematy podpisów do zastosowania przy dodawaniu zrzutów ekranu:

\begin{itemize}
  \item \textit{Główne okno aplikacji (PyQt6) podczas wsadowej analizy filmów; lewy panel zawiera tabelę z~wynikami detekcji watermarków (kolumny: plik, typ, status, \% AI, C2PA, CSV), prawy panel prezentuje aktualnie analizowaną klatkę wideo oraz powiększony obszar wykrycia.}
  \item \textit{Katalog raportu wygenerowanego przez aplikację dla jednego pliku wideo; widoczne są zapisane klatki z~zaznaczonymi wykryciami OCR, obraz mediany temporalnej, wizualizacja FFT i~plik \texttt{report.csv} z~metadanymi detekcji.}
  \item \textit{Konsola logów aplikacji prezentująca kolejne etapy analizy: inicjalizację silnika OCR (EasyOCR/PaddleOCR), próbkowanie klatek z~zadanym \textit{sample rate}, wykrycie regionu tekstowego z~podanym poziomem pewności oraz zapis raportu do~katalogu wynikowego.}
  \item \textit{Wynik analizy pliku \texttt{000000passport\_fake.mp4}; w~tabeli status \texttt{🔴 AI DETECTED} z~typem detekcji, w~kolumnie C2PA brak metadanych Content Credentials; panel Zoom prezentuje powiększony obszar wykrycia.}
\end{itemize}

\section{Źródła i inspiracje techniczne}

Rozwój systemu był wspierany przez dokumentacje techniczne bibliotek \textbf{OpenCV}, \textbf{PyQt6}, \textbf{EasyOCR} i~\textbf{PaddleOCR}, a~w~warstwie koncepcyjnej przez publiczne materiały dotyczące standardu \textbf{C2PA} (\url{https://c2pa.org}), dokumentację systemu \textbf{SynthID} (Google DeepMind) oraz publikacje i~zasoby omówione w~rozdziałach~2--3 niniejszej pracy.

Doskonalenie detekcji miało charakter iteracyjny. Wyniki własnego detektora były porównywane z~zachowaniem publicznie dostępnych generatorów wideo (Runway, Sora, Luma Dream Machine) i~komercyjnych detektorów, analizowane były materiały demonstracyjne publikowane przez dostawców modeli, a~skuteczność poszczególnych heurystyk weryfikowano po~typowych operacjach na~plikach: kompresji, skalowaniu, rekodowaniu i~konwersji formatu. Dzięki temu projekt ewoluował w~kierunku wielosygnałowego detektora, który łączy analizę OCR, przepływu optycznego, artefaktów FFT i~metadanych pochodzenia.

\section{Ewaluacja jakościowa}

Ewaluacja systemu w~niniejszej pracy ma charakter jakościowy i~eksploracyjny. Oceniano, czy aplikacja potrafi wykrywać charakterystyczne sygnały obecne w~materiałach z~popularnych generatorów wideo, czy wyniki są stabilne po~kompresji oraz czy system unika fałszywych alarmów dla materiałów zawierających legalne nakładki graficzne (billboardy, dolne napisy telewizyjne, transmisje strumieniowe).

Istotną cechą zaproponowanego systemu jest dostarczanie wyników pośrednich, co pozwala sprawdzić, czy detekcja wynikała z~wykrycia tekstu przez OCR, z~utrzymywania się statycznej nakładki w~kolejnych klatkach, z~anomalii w~analizie FFT czy z~wykrycia metadanych C2PA. Zwiększa to wiarygodność systemu i~stanowi przewagę nad narzędziami zwracającymi wyłącznie pojedynczy procentowy wskaźnik bez uzasadnienia --- jak pokazano w~rozdziale~4 na~przykładzie \textit{detectvideo.ai}, który dla~80\% badanych filmów zwrócił wynik~55\% (\textit{inconclusive}), nie podejmując żadnej decyzji klasyfikacyjnej.

\section{Wnioski implementacyjne}

Przeprowadzona implementacja potwierdziła, że skuteczny detektor materiałów AI nie powinien opierać się na~jednym sygnale. Znacznie lepsze rezultaty daje połączenie kilku niezależnych źródeł informacji: OCR, analizy ruchu, artefaktów częstotliwościowych, watermarkingu niewidzialnego i~metadanych proweniencji. Doświadczenia z~budowy aplikacji pokazują jednocześnie, że równie ważna jak warstwa algorytmiczna jest jakość interfejsu, czytelność raportowania i~dostępność artefaktów diagnostycznych --- szczególnie przy analizie trudnych przypadków, w~których automatyczna decyzja wymaga ludzkiej weryfikacji.
